% Chapter 13 back matter: physical PDF pp. 585--586 (printed pp. 562--563).
% Inventory: five problems and ten references; no figures or footnotes.

\chapterbackmatter{Problems}

\begin{enumerate}
\item Consider the process $e^+ + e^- \to \pi^+ + \pi^-$ in the center-of-mass
frame at an energy of $1$ GeV and scattering angle of $90^\circ$. Suppose that
by measuring the energies of the final pions we determine that an energy of not
more than $E_T \ll 1$ GeV is emitted in soft photons. How does the reaction rate
depend on $E_T$?

\item Consider a massless spinless particle described by a scalar field $\phi$,
whose interaction Lagrangian density is of the form $\phi(x)J(x)$, where $J(x)$
involves only massive particle fields. Derive a formula for the rate of emission
of arbitrary numbers of soft scalar particles in a process $\alpha\to\beta$,
with these soft scalars having a total energy less than some small quantity
$E_T$. Include radiative corrections due to soft scalars with energy less than
some small quantity $\Lambda$.

\item Derive a formula for the next term beyond Eq.
\eqref{eq:13.5.14} for
low-energy photon scattering on an arbitrary target.

\item Suppose that a spin one particle of very small mass $m$ is described by a
vector field $V^\mu(x)$, which couples only to a much heavier fermion described
by a Dirac field $\psi(x)$, with an interaction Lagrangian density of the form
$g V^\mu \bar\psi\gamma_\mu\psi$. Suppose that the heavy fermion normally decays
into other particles that have no interaction with $V^\mu$, releasing an energy
$W$ much greater than $m$. Consider such a decay process, but where an
additional $V^\mu$ particle is emitted along with the other decay products,
with the vector particle energy less than an upper limit $E$, in the range
$W\gg E\gg m$. How does the rate of this process depend on $E$ and $m$? Ignore
radiative corrections.

\item Prove that Eq. \eqref{eq:13.6.8} holds when
$p\cdot(q_1+\cdots+q_N)=0$.
\end{enumerate}

\chapterbackmatter{References}

\begin{enumerate}
\item \label{ref:13.1}F. Bloch and A. Nordsieck, \emph{Phys. Rev.} \textbf{37},
54 (1937); D. R. Yennie, S. C. Frautschi, and H. Suura, \emph{Ann. Phys. (NY)}
\textbf{13}, 379 (1955).

\item \label{ref:13.2}S. Weinberg, \emph{Phys. Lett.} \textbf{9}, 357 (1964);
\emph{Phys. Rev.} \textbf{135}, B1049 (1964).

\item \label{ref:13.3}S. Weinberg, \emph{Phys. Rev.} \textbf{140}, B515
(1965).

\item \label{ref:13.4}This phase was encountered in the perturbation series for
non-relativistic Coulomb scattering, by R. H. Dalitz, \emph{Proc. Roy. Soc.
London} \textbf{206}, 509 (1951).

\item \label{ref:13.5}See, e.g., L. I. Schiff, \emph{Quantum Mechanics}
(McGraw-Hill, New York, 1949): Section 20.

\item \label{ref:13.6}T. D. Lee and M. Nauenberg, \emph{Phys. Rev.}
\textbf{133}, B1549 (1964). Also see T. Kinoshita, \emph{J. Math. Phys.}
\textbf{3}, 650 (1962); G. Sterman and S. Weinberg, \emph{Phys. Rev. Lett.}
\textbf{39}, 1416 (1977).

\item \label{ref:13.7}G. Sterman and S. Weinberg,
Ref.~\hyperref[ref:13.6]{6}.

\item \label{ref:13.8}F. E. Low, \emph{Phys. Rev.} \textbf{96}, 1428 (1954);
M. Gell-Mann and M. L. Goldberger, \emph{Phys. Rev.} \textbf{96}, 1433 (1954).
Also see S. Weinberg, in \emph{Lectures on Elementary Particles and Quantum
Field Theory---1970 Brandeis Summer Institute in Theoretical Physics}, ed. by
S. Deser, M. Grisaru, and H. Pendleton (MIT Press, Cambridge, MA, 1970).

\item \label{ref:13.9}E. J. Williams, \emph{Kgl. Dan. Vid. Sel. Mat.-fys. Medd.}
\textbf{XIII}, No. 4 (1935).

\item \label{ref:13.10}H. A. Bethe and E. E. Salpeter, \emph{Phys. Rev.}
\textbf{82}, 309 (1951); \textbf{84}, 1232 (1951).
\end{enumerate}
